\section{Background and related work}
\label{sec:bg}

We sketch prior work along the two axes that frame the empty cell
\sysname enters---\textbf{capability mediation} and \textbf{userspace
IPC}---and defer the longer discussion to \Cref{sec:rw}.

\subsection{Capability systems}
\label{sec:bg:cap}

Capabilities go back to Dennis and Van~Horn~\cite{dennis1966}.
Redell's thesis~\cite{redell74} supplied the piece that matters here:
the epoch-versioned credential, the ancestor of every revocable
capability since. KeyKOS~\cite{keykos} and EROS~\cite{shapiro1999eros}
showed the pattern working at operating-system scale,
CapROS~\cite{capros} carried it forward, and seL4~\cite{klein2009sel4}
closed the loop with a full functional-correctness proof of a
capability microkernel.

Two recent systems form the immediate peer set.
\textbf{Cornucopia\,Reloaded}~\cite{cornucopia-reloaded} (ASPLOS~'24)
brings epoch revocation to a CHERI hardware substrate with $O(1)$
bulk invalidation; our revocation scheme \emph{is} theirs, ported to
software and held inside the Linux process boundary. We say this
plainly because the algorithmic credit belongs to them---what we add
is the realisation without CHERI.
\textbf{HongMeng}~\cite{hongmeng} (OSDI~'24) is the closest
competitor on the Linux-compatible, capability-based, fast-IPC
triple. It is a kernel underneath Linux rather than a runtime above
it, and its authors report retreating from ``pure'' capability chains
to ``address tokens'' because chain-revocation cost hurt the common
case. We sit on the other side of the kernel boundary, and
\Cref{sec:disc:hongmeng} discusses what that buys and costs.

\textbf{CAP-VMs}~\cite{capvms} (OSDI~'22) pairs capabilities with
fast IPC in userspace, but requires CHERI/Morello. Our positioning
relative to it is simple: the same conceptual layer, on stock Linux.

\subsection{Userspace IPC primitives}
\label{sec:bg:ipc}

\textbf{Aeron}~\cite{aeron-real-logic} runs single-producer,
single-consumer IPC over an mmap'd ring buffer at a published
$\sim$250\,ns round trip for 100-byte payloads on commodity Xeon
hardware. That is the floor we measure against. ByteDance's
shmipc~\cite{bytedance-shmipc} is architecturally similar and
production-deployed.

\textbf{io\_uring}~\cite{iouring} put the same template---mmap'd
SQ/CQ rings with futex-style notification---inside the Linux kernel.
Our ring buffer reuses the template one layer up: the kernel never
sees our messages, and we never pay the \texttt{enter} syscall.

\textbf{eRPC}~\cite{erpc} (NSDI~'19) holds the lowest peer-reviewed
RPC latency we know of ($\sim$2.3\,$\mu$s for a single-hop
request/response). It is an RPC stack rather than a raw transport, so
we cite it as the network-side reference point rather than fighting
it on raw IPC RTT.

\textbf{LITESHIELD}~\cite{liteshield} (USENIX ATC~'25) is the
architectural neighbour: a userspace microkernel-style runtime on
Linux with shared-memory IPC and a host interface of roughly 22
syscalls---and no capability layer. We cite it prominently because
the structural overlap is real; \sysname is, in effect, the
capability layer over that shape of system.

\subsection{The empty cell}

\Cref{tab:landscape} lays out the grid. The intersection
\textit{Linux-compatible} $\cap$ \textit{no CHERI} $\cap$
\textit{capability-mediated IPC} $\cap$ \textit{measured against
Aeron-class numbers} contains, as far as we can tell, nothing.
\sysname is our attempt to put something there.

\begin{table}[t]
\centering\footnotesize
\setlength{\tabcolsep}{4pt}
\begin{tabular}{lccccc}
\toprule
System            & Linux & no CHERI & Caps & IPC fastpath & Userspace \\
\midrule
seL4              & no  & yes & yes & yes & no  \\
HongMeng          & yes & yes & weak & yes & no  \\
gVisor            & yes & yes & no  & no  & yes \\
Firecracker       & yes & yes & no  & no  & yes \\
CAP-VMs           & yes & no  & yes & yes & yes \\
LITESHIELD        & yes & yes & no  & yes & yes \\
Aeron             & yes & yes & no  & yes & yes \\
io\_uring         & yes & yes & no  & yes & no  \\[2pt]
\textbf{\sysname} & \textbf{yes} & \textbf{yes} & \textbf{yes} & \textbf{yes} & \textbf{yes} \\
\bottomrule
\end{tabular}
\caption{The cell \sysname enters. ``weak'' marks HongMeng's own
retreat from capability chains to address tokens.}
\label{tab:landscape}
\end{table}
